% GENERATED FILE. Edit canonical lessons, then run npm run generate:books.
% canonical-source-hash: fnv1a64:ae59b14ddbf5c165
% canonical-lessons: FR-C219-ennuyeux-ennuyeuse, FR-C219-amusant-amusante, FR-C219-etrange, FR-C219-different-differente, FR-C219-pareil-pareille

\chapter{Boring, Fun, Strange, Different, Same}
\label{ch:fr-boring-fun-strange-different-same}
\hlchaptermodality{219}

\begin{chapteropening}
\textbf{By the end of this chapter:} \emph{I can say boring, fun, strange, different and the same.}

Complete the last lesson of chapter 219: I can say boring, fun, strange, different and the same.
\end{chapteropening}

\section[ennuyeux / ennuyeuse]{ennuyeux / ennuyeuse — boring}

\label{lesson:FR-C219-ennuyeux-ennuyeuse}



\begin{quote}
Before the new one: say the French for nervous, then the French for quiet.
\end{quote}

\subsection*{You'll want to know: ennuyeux / ennuyeuse}
\textbf{ennuyeux / ennuyeuse} — \textquotedblleft{}boring\textquotedblright{}.

\begin{cousinweb}
\textbf{ennuyeux} — built from the noun beside \emph{s'ennuyer}, with \emph{-eux}, \textquotedblleft{}full of\textquotedblright{}: full of boredom.
\end{cousinweb}

\begin{grammarlens}[title={What you've built}]
One more word for this chapter.
\end{grammarlens}

\subsection*{Your turn}
\begin{itemize}
\raggedright
  \item \emph{Say it:} \emph{ennuyeux / ennuyeuse}
  \item \emph{Say it:} \emph{ennuyeux / ennuyeuse}, once more
  \item \emph{Say it:} say \emph{ennuyeux / ennuyeuse}
  \item \emph{Recall:} say the French for nervous, then the French for quiet, then say \emph{ennuyeux / ennuyeuse} again
\end{itemize}

\begin{tcolorbox}[breakable,colback=teal!4,colframe=teal!35!black,title={Before you move on}]
What does ennuyeux / ennuyeuse mean? (\textquotedblleft{}Boring\textquotedblright{}.) Say it once more.
\end{tcolorbox}
\section[amusant / amusante]{amusant / amusante — fun, amusing}

\label{lesson:FR-C219-amusant-amusante}



\begin{quote}
Before the new one: say the French for quiet, then the French for boring.
\end{quote}

\subsection*{You'll want to know: amusant / amusante}
\textbf{amusant / amusante} — \textquotedblleft{}fun, amusing\textquotedblright{}.

\begin{cousinweb}
\textbf{amusant} is the present participle of the verb in \emph{s'amuser}, \textquotedblleft{}to amuse\textquotedblright{}: amusing. \emph{Amuser} $\leftarrow$ Old French \emph{amuser}, \textquotedblleft{}to keep idle\textquotedblright{}, from \emph{muser}, \textquotedblleft{}to gape\textquotedblright{} (EtymDB 2.1).
\end{cousinweb}

\begin{grammarlens}[title={What you've built}]
One more word for this chapter.
\end{grammarlens}

\subsection*{Your turn}
\begin{itemize}
\raggedright
  \item \emph{Say it:} \emph{amusant / amusante}
  \item \emph{Say it:} \emph{amusant / amusante}, once more
  \item \emph{Say it:} say \emph{amusant / amusante}
  \item \emph{Recall:} say the French for quiet, then the French for boring, then say \emph{amusant / amusante} again
\end{itemize}

\begin{tcolorbox}[breakable,colback=teal!4,colframe=teal!35!black,title={Before you move on}]
What does amusant / amusante mean? (\textquotedblleft{}Fun, amusing\textquotedblright{}.) Say it once more.
\end{tcolorbox}
\section[étrange]{étrange — strange}

\label{lesson:FR-C219-etrange}



\begin{quote}
Before the new one: say the French for boring, then the French for fun.
\end{quote}

\subsection*{You'll want to know: étrange}
\textbf{étrange} — \textquotedblleft{}strange\textquotedblright{}.

\begin{cousinweb}
\textbf{étrange} $\leftarrow$ Old French \emph{estrange} $\leftarrow$ Latin \emph{extraneus}, \textquotedblleft{}from outside, foreign\textquotedblright{} — inherited. (EtymDB 2.1)
\end{cousinweb}

\begin{grammarlens}[title={What you've built}]
One more word for this chapter.
\end{grammarlens}

\subsection*{Your turn}
\begin{itemize}
\raggedright
  \item \emph{Say it:} \emph{étrange}
  \item \emph{Say it:} \emph{étrange}, once more
  \item \emph{Say it:} say \emph{étrange}
  \item \emph{Recall:} say the French for boring, then the French for fun, then say \emph{étrange} again
\end{itemize}

\begin{tcolorbox}[breakable,colback=teal!4,colframe=teal!35!black,title={Before you move on}]
What does étrange mean? (\textquotedblleft{}Strange\textquotedblright{}.) Say it once more.
\end{tcolorbox}
\section[différent / différente]{différent / différente — different}

\label{lesson:FR-C219-different-differente}



\begin{quote}
Before the new one: say the French for fun, then the French for strange.
\end{quote}

\subsection*{You'll want to know: différent / différente}
\textbf{différent / différente} — \textquotedblleft{}different\textquotedblright{}.

\begin{cousinweb}
\textbf{différent} $\leftarrow$ Latin \emph{differens}, \textquotedblleft{}carrying apart\textquotedblright{}, from \emph{differre} (Lewis \& Short) — a learned borrowing.
\end{cousinweb}

\begin{grammarlens}[title={What you've built}]
One more word for this chapter.
\end{grammarlens}

\subsection*{Your turn}
\begin{itemize}
\raggedright
  \item \emph{Say it:} \emph{différent / différente}
  \item \emph{Say it:} \emph{différent / différente}, once more
  \item \emph{Say it:} say \emph{différent / différente}
  \item \emph{Recall:} say the French for fun, then the French for strange, then say \emph{différent / différente} again
\end{itemize}

\begin{tcolorbox}[breakable,colback=teal!4,colframe=teal!35!black,title={Before you move on}]
What does différent / différente mean? (\textquotedblleft{}Different\textquotedblright{}.) Say it once more.
\end{tcolorbox}
\section[pareil / pareille]{pareil / pareille — the same}

\label{lesson:FR-C219-pareil-pareille}



\begin{quote}
Before the new one: say the French for strange, then the French for different.
\end{quote}

\subsection*{You'll want to know: pareil / pareille}
\textbf{pareil / pareille} — \textquotedblleft{}the same\textquotedblright{}.

\begin{cousinweb}
\textbf{pareil} $\leftarrow$ Middle French \emph{pareil} (EtymDB 2.1) $\leftarrow$ Vulgar Latin \emph{*pariculus}, from Latin \emph{par}, \textquotedblleft{}equal\textquotedblright{} (Lewis \& Short) — inherited.
\end{cousinweb}

\begin{grammarlens}[title={What you've built}]
One more word for this chapter.
\end{grammarlens}

\subsection*{Your turn}
\begin{itemize}
\raggedright
  \item \emph{Say it:} \emph{pareil / pareille}
  \item \emph{Say it:} \emph{pareil / pareille}, once more
  \item \emph{Say it:} say \emph{pareil / pareille}
  \item \emph{Recall:} say the French for strange, then the French for different, then say \emph{pareil / pareille} again
\end{itemize}

\begin{tcolorbox}[breakable,colback=teal!4,colframe=teal!35!black,title={Before you move on}]
What does pareil / pareille mean? (\textquotedblleft{}The same\textquotedblright{}.) Say it once more.
\end{tcolorbox}
